\documentclass[11pt]{article}

\usepackage[margin=1in]{geometry}
\usepackage[T1]{fontenc}
\usepackage{lmodern}

\usepackage{amsmath, amssymb}

\usepackage{setspace}

\usepackage[most]{tcolorbox}
\usepackage{xcolor}

\tcbset{
  enhanced,
  boxrule=0pt,
  frame hidden,
  breakable
}

\usepackage{fancyhdr}
\pagestyle{fancy}
\fancyhf{}
\setlength{\headheight}{52pt}
\renewcommand{\headrulewidth}{0pt}

\newcommand{\Course}{Math 4610 Spring 2026}
\newcommand{\Instructor}{Discussion 4}
\newcommand{\AssignmentType}{Discussion} % or Discussion
\newcommand{\AssignmentNumber}{4}

\newcommand{\Contributors}{
Marks, Stephen\\
}

\fancyhead[L]{\Course\\ \Instructor}
\fancyhead[R]{\Contributors}

\newcommand{\ProblemDivider}[1]{%
  \vspace{1.0em}
  \noindent\textbf{#1}\par
  \vspace{0.35em}
}

\definecolor{problemBack}{RGB}{235,242,250}
\definecolor{problemFrame}{RGB}{70,110,160}
\definecolor{exerciseGray}{gray}{0.96}
\definecolor{exerciseGrayFrame}{gray}{0.45}

\newtcolorbox{problemBox}{
  colback=exerciseGray,
  colframe=exerciseGrayFrame,
  arc=2mm,
  left=8pt,right=8pt,top=8pt,bottom=8pt
}

\newtcolorbox{solutionBox}{
  colback=white,
  colframe=white,
  arc=0mm,
  left=0pt,right=0pt,top=0pt,bottom=0pt
}

\newcommand{\ProblemAnnounce}{\noindent}
\newcommand{\SolutionAnnounce}{\noindent\textbf{Solution.}\par\medskip}

\newcommand{\Exercise}[1]{\textbf{Exercise~#1.}\;}

\newcommand{\R}{\mathbb{R}}
\newcommand{\C}{\mathbb{C}}
\begin{document}

\begin{center}
  {\Large \textbf{\AssignmentType~\AssignmentNumber}}
\end{center}
\vspace{0.75em}

\ProblemDivider{1}

\begin{problemBox}
\ProblemAnnounce
Suppose $T\in \mathcal{L}(\mathbb{F}^3)$ and
\[
M(T)=
\begin{pmatrix}
2 & 0 & 1\\
0 & 1 & -1\\
0 & 0 & 1
\end{pmatrix},
\]
with respect to the standard basis $e_1,e_2,e_3$. Find the minimal polynomial of $T$.
\end{problemBox}

\begin{solutionBox}
\SolutionAnnounce
\doublespacing

Note that $M(T)$ is upper triangular, so the diagonal is comprised of eigenvalues
$\lambda_1=2$ and $\lambda_2=\lambda_3=1$.
Thus $(T-2I)(T-I)^2=0$ by 5.40 \& 5.41.

Compute
\[
T-I=
\begin{pmatrix}
2 & 0 & 1\\
0 & 1 & -1\\
0 & 0 & 1
\end{pmatrix}
-
\begin{pmatrix}
1 & 0 & 0\\
0 & 1 & 0\\
0 & 0 & 1
\end{pmatrix}
=
\begin{pmatrix}
1 & 0 & 1\\
0 & 0 & -1\\
0 & 0 & 0
\end{pmatrix}
\neq 0.
\]

And,
\[
(T-I)^2 =
\begin{pmatrix}
1 & 0 & 1\\
0 & 0 & -1\\
0 & 0 & 0
\end{pmatrix}
\begin{pmatrix}
1 & 0 & 1\\
0 & 0 & -1\\
0 & 0 & 0
\end{pmatrix}
=
\begin{pmatrix}
1 & 0 & 1\\
0 & 0 & 0\\
0 & 0 & 0
\end{pmatrix}
\neq 0.
\]

And,
\[
(T-2I)=
\begin{pmatrix}
0 & 0 & 1\\
0 & -1 & -1\\
0 & 0 & -1
\end{pmatrix}
\neq 0.
\]

Since $(T-2I)(T-I)\neq 0$, $(T-2I)\neq 0$, $(T-I)^2\neq 0$, and $(T-I)\neq 0$,
no factor annihilates $T$. Therefore, the minimal polynomial of $T$ must be
\[
f(z)=(z-2)(z-1)^2,
\]
which is supported by 5.44.

\end{solutionBox}
\newpage

\ProblemDivider{2}

\begin{problemBox}
\ProblemAnnounce
Suppose $T\in \mathcal{L}(\mathbb{F}^3)$ and
\[
M(T)=
\begin{pmatrix}
2 & 0 & 0\\
0 & 1 & 0\\
0 & 0 & 1
\end{pmatrix},
\]
with respect to the standard basis $e_1,e_2,e_3$. Find the minimal polynomial of $T$.
\end{problemBox}

\begin{solutionBox}
\SolutionAnnounce
\doublespacing

Note that $M(T)$ is diagonal, thus diagonalizable and upper triangular. So the diagonal is comprised of
eigenvalues $\lambda_1=2$ and $\lambda_2=\lambda_3=1$.
Thus $(T-2I)(T-I)^2=0$. However, the minimal polynomial is at most $(z-2)(z-1)$
because the factors of the minimal polynomial, due to diagonalizability, must be linear and distinct by 5.62.

Compute
\[
T-I=
\begin{pmatrix}
1 & 0 & 0\\
0 & 0 & 0\\
0 & 0 & 0
\end{pmatrix}
\neq 0,
\quad
T-2I=
\begin{pmatrix}
0 & 0 & 0\\
0 & -1 & 0\\
0 & 0 & -1
\end{pmatrix}
\neq 0.
\]

Since $(T-2I)\neq 0$ and $(T-I)\neq 0$, no factor annihilates $T$.
Therefore, the minimal polynomial is
\[
f(z)=(z-2)(z-1).
\]

\end{solutionBox}

\end{document}

